% ============================================================
% 02-related-work.tex
% ============================================================
\section{Related Work}
\label{sec:related}

This section reviews established retrieval and reasoning paradigms in the literature, examining their operational mechanisms and limitations in institutional counseling.

%------------------------------------------------
\subsection{Vietnamese NLP Foundations and Linguistic Specifics}
\label{sec:vietnamese_nlp}

Unlike whitespace-delimited languages, Vietnamese is an isolating, tonal language where white spaces separate syllables rather than words~\cite{nguyen2020phobert}, with over 85\% of lexical units forming multi-syllabic compound words (e.g., \textit{``h\d{o}c ph\'i''} [tuition fee]). While domain-adapted bi-encoders like PhoBERT capture semantic paraphrasing, Vietnamese higher education decrees rely heavily on rigid Sino-Vietnamese legal terminology and alphanumeric cohort codes. This dual characteristic motivates coupling syllable-aware dense embeddings with exact lexical token matching to avoid both semantic drift and alphanumeric recall failures.

%------------------------------------------------
\subsection{Sparse and Dense Retrieval}
\label{sec:sparse_dense}

\textbf{Lexical Retrieval (BM25):} Sparse keyword search using BM25~\cite{robertson2009bm25} scores documents via term frequency ($\mathrm{tf}$) and inverse document frequency ($\mathrm{IDF}$) over an inverted index ($k_1 = 1.5, b = 0.75$). BM25 guarantees exact lexical matching for rigid alphanumeric identifiers (decision codes, cohort labels, course codes), but cannot capture semantic paraphrases.

\textbf{Dense Semantic Retrieval:} Dense retrieval systems~\cite{karpukhin2020dpr,lewis2020rag} map queries and passages into embedding spaces using bi-encoders, computing relevance via vector similarity. While dense models capture paraphrastic intent, rare alphanumeric identifiers and out-of-vocabulary regulatory tokens can remain difficult and may introduce cross-cohort noise, motivating hybrid retrieval.

%------------------------------------------------
\subsection{Hybrid Retrieval and Reciprocal Rank Fusion}
\label{sec:hybrid_rrf}

To merge BM25 and dense candidates without score calibration, Reciprocal Rank Fusion (RRF)~\cite{cormack2009rrf} sorts documents by reciprocal position: $\RRF(d) = \sum_{s \in \mathcal{M}} \frac{1}{k + r_s(d)}$ ($k = 60$), while neural cross-encoders~\cite{nogueira2019bert,xiao2023bge} jointly encode query--passage pairs for reranking. Adaptive-RAG further demonstrates that retrieval strategies can be selected according to question complexity~\cite{jeong2024adaptiverag}. These approaches motivate dynamic retrieval, but do not directly allocate institution-specific relational, tabular, and narrative knowledge to isolated specialist paths.

%------------------------------------------------
\subsection{Knowledge Graphs and GraphRAG}
\label{sec:kg_graphrag}

Knowledge graphs represent domain knowledge as structured triples $(h, r, t)$~\cite{pan2024unifying}. GraphRAG~\cite{edge2024graphrag} leverages graph topology to uncover multi-hop associative paths escaping flat chunking, benefiting curriculum structures and prerequisite chains. However, universal graph extraction across institutional corpora produces hallucinated relationships, while forcing narrative legal prose into graphs is structurally unnatural, motivating selective rather than universal graph use.

%------------------------------------------------
\subsection{Agentic RAG and Multi-Agent Orchestration}
\label{sec:agentic_multiagent}

ReAct demonstrates interleaved reasoning and action, AutoGen provides a framework for multi-agent conversation, and Toolformer studies when and how language models invoke external APIs~\cite{yao2023react,wu2023autogen,schick2023toolformer}. For institutional counseling, these capabilities raise design risks when applied without explicit constraints: broad tool exposure can increase tool and parameter errors; peer-to-peer exchanges can increase coordination cost; a shared index can mix evidence across domains; and a common retrieval path can obscure opportunities for direct graph traversal or structured lookup. These risks motivate the bounded, supervisor-routed specialization evaluated in this work.

%------------------------------------------------
\subsection{Educational and University Counseling Systems}
\label{sec:educational_counseling}

Educational chatbots assist in advising and learning support~\cite{wollny2021chatbots}, with systems like REBot~\cite{ma2025rebot} combining text retrieval with category-guided GraphRAG for regulations. \system{} departs from prior advisory systems across three dimensions: (1)~\emph{Task Scope}: while prior systems focus on single-turn text-QA over general regulations, \system{} models 113 curricula as relational property graphs and incorporates deterministic tools (\texttt{tinh\_toan\_hoc\_phi}, \texttt{tinh\_tien\_hoc\_bong}) for conditional tuition schedules; (2)~\emph{Multi-Agent Routing}: rather than monolithic pipelines, a LangGraph supervisor coordinates four bounded specialists with isolated evidence spaces and asymmetric tools; and (3)~\emph{Evaluation Rigor}: evaluating end-to-end multi-agent faithfulness via Ragas and programmatic tool execution rather than shallow classification accuracy.

%------------------------------------------------
\subsection{Comparative Synthesis and Research Gaps}
\label{sec:gaps}

Synthesizing established paradigms highlights clear trade-offs across institutional counseling workflows~\cite{cormack2009rrf,edge2024graphrag,jeong2024adaptiverag,karpukhin2020dpr,robertson2009bm25,schick2023toolformer,wu2023autogen,yao2023react}: lexical matching ensures exact identifier recall but misses paraphrases; dense embeddings capture semantics but cause cross-cohort confusion; GraphRAG models relations but degrades on unstructured text; and unconstrained multi-agent frameworks lack execution determinism. These trade-offs crystallize three core research gaps:

\begin{itemize}
    \setlength{\itemsep}{1pt}
    \setlength{\parsep}{0pt}
    \item \textbf{Gap 1 (Lack of Bounded Domain-Specialized Agent Orchestration):} Prior frameworks rely on monolithic pipelines or unconstrained agent meshes without execution boundaries, increasing the risks of tool mis-selection, parameter hallucination, and high coordination costs.
    \item \textbf{Gap 2 (The Uniform Representation Bottleneck):} Forcing relational curricula, rate matrices, and regulatory prose into a uniform retrieval index fragments prerequisite chains and induces relation hallucinations across structurally incompatible evidence spaces.
    \item \textbf{Gap 3 (Limited Evidence for Representation-Aware Retrieval):} Existing systems provide individual lexical, dense, hybrid, and graph retrieval mechanisms, but evidence remains limited on how an integrated representation-aware retrieval stack performs overall and across heterogeneous institutional domains.
\end{itemize}

These limitations motivate our centralized supervisor-routed architecture and heterogeneous knowledge allocation detailed in Section~\ref{sec:model}.
